
\documentclass{article}
\usepackage{colortbl}
\usepackage{makecell}
\usepackage{multirow}
\usepackage{supertabular}

\begin{document}

\newcounter{utterance}

\centering \large Interaction Transcript for game `eqbench', experiment `eq\_bench', episode 166 with reflexion{-}llama{-}3.1{-}8b{-}instruct.
\vspace{24pt}

{ \footnotesize  \setcounter{utterance}{1}
\setlength{\tabcolsep}{0pt}
\begin{supertabular}{c@{$\;$}|p{.15\linewidth}@{}p{.15\linewidth}p{.15\linewidth}p{.15\linewidth}p{.15\linewidth}p{.15\linewidth}}
    \# & \multicolumn{2}{c}{Player} && \multicolumn{2}{c}{Game Master} \\
    \hline

    \theutterance \stepcounter{utterance}  
    & & & \multicolumn{4}{p{0.6\linewidth}}{
        \cellcolor[rgb]{0.9,0.9,0.9}{
            \makecell[{{p{\linewidth}}}]{
                \texttt{\tiny{[P1$\langle$GM]}}
                \\ 
\texttt{Your task is to predict the likely emotional responses of a character in this dialogue:} \\
\\ 
\texttt{Jonny: You're not like the others, are you, Eva?} \\
\texttt{Eva: And what makes you say that?} \\
\texttt{Jonny: Well, for starters, you can see me. And you don't scream and run away.} \\
\texttt{Eva: (laughs) I suppose that does make me a bit different.} \\
\texttt{Jonny: More than a bit, I'd say. } \\
\texttt{Eva: I've always been... sensitive to things others can't see. } \\
\texttt{Jonny: Like a sixth sense?} \\
\texttt{Eva: Something like that. } \\
\texttt{Jonny: Must be lonely, seeing things no one else can.} \\
\texttt{Eva: It can be... but it also allows me to meet interesting individuals. Like you.} \\
\texttt{Jonny: (laughs) I'm flattered. Not many people find a ghost interesting.} \\
\texttt{Eva: You're more than just a ghost, Jonny. You're a person, just... not in the traditional sense.} \\
\texttt{Jonny: That's... unexpectedly kind of you to say, Eva.} \\
\texttt{Eva: (smiles) Well, it's true.} \\
\texttt{Jonny: (pauses) Do you... Do you ever get scared? Dealing with... well, things like me?} \\
\texttt{Eva: Sometimes. But not with you. You're different.} \\
\texttt{Jonny: Different how?} \\
\texttt{Eva: You're... (pauses) You're not like the others.} \\
\texttt{Jonny: (laughs) Now who's sounding like a broken record?} \\
\texttt{Eva: (laughs) I suppose I walked into that one.} \\
\texttt{Jonny: (softly) Eva... } \\
\texttt{Eva: Yes, Jonny?} \\
\texttt{Jonny: Never mind. It's... it's nothing.} \\
\texttt{{[}End dialogue{]}} \\
\\ 
\texttt{At the end of this dialogue, Eva would feel...} \\
\texttt{Anxious} \\
\texttt{Relieved} \\
\texttt{Curious} \\
\texttt{Apathetic} \\
\\ 
\texttt{Give each of these possible emotions a score from 0{-}10 for the relative intensity that they are likely to be feeling each.} \\
\\ 
\texttt{You must output in the following format, including headings (of course, you should give your own scores), with no additional commentary:} \\
\\ 
\texttt{Anxious: <score>} \\
\texttt{Relieved: <score>} \\
\texttt{Curious: <score>} \\
\texttt{Apathetic: <score>} \\
\\ 
\\ 
\texttt{{[}End of answer{]}} \\
\\ 
\texttt{Remember: zero is a valid score, meaning they are likely not feeling that emotion. You must score at least one emotion > 0.} \\
\\ 
\texttt{Your answer:} \\
            }
        }
    }
     \\ \\

    \theutterance \stepcounter{utterance}  
    & \multicolumn{4}{p{0.6\linewidth}}{
        \cellcolor[rgb]{0.561,0.737,0.561}{
            \makecell[{{p{\linewidth}}}]{
                \texttt{\tiny{[P1$\rangle$GM]}}
                \texttt{Anxious: 6} \\
\texttt{Relieved: 8} \\
\texttt{Curious: 0} \\
\texttt{Apathetic: 0} \\
            }
        }
    }
    & & \\ \\

    \theutterance \stepcounter{utterance}  
    & & & \multicolumn{2}{p{0.3\linewidth}}{
        \cellcolor[rgb]{0.9,0.9,0.9}{
            \makecell[{{p{\linewidth}}}]{
                \texttt{\tiny{[GM$|$GM]}}
                \texttt{Anxious: 6} \\
\texttt{Relieved: 8} \\
\texttt{Curious: 0} \\
\texttt{Apathetic: 0} \\
            }
        }
    }
    & & \\ \\

    \theutterance \stepcounter{utterance}  
    & & & \multicolumn{2}{p{0.3\linewidth}}{
        \cellcolor[rgb]{0.9,0.9,0.9}{
            \makecell[{{p{\linewidth}}}]{
                \texttt{\tiny{[GM$|$GM]}}
                \texttt{Anxious: 2} \\
\texttt{Relieved: 0} \\
\texttt{Curious: 8} \\
\texttt{Apathetic: 0} \\
            }
        }
    }
    & & \\ \\

    \theutterance \stepcounter{utterance}  
    & & & \multicolumn{2}{p{0.3\linewidth}}{
        \cellcolor[rgb]{0.9,0.9,0.9}{
            \makecell[{{p{\linewidth}}}]{
                \texttt{\tiny{[GM$|$GM]}}
                \texttt{game\_result = LOSE} \\
            }
        }
    }
    & & \\ \\

    \theutterance \stepcounter{utterance}  
    & & \multicolumn{4}{p{0.6\linewidth}}{
        \cellcolor[rgb]{0.95,0.95,0.95}{
            \makecell[{{p{\linewidth}}}]{
                \texttt{\tiny{[P1$\leftrightarrow$P]}}
                \texttt{Post{-}episode Reflexion (not sent to GM):} \\
\texttt{New plan: } \\
\\ 
\texttt{Upon re{-}reading the dialogue, I realized that I focused too much on Eva's responses to Jonny's questions and didn't consider the emotional context of the conversation. I should have paid more attention to the subtle cues and undertones in Eva's words and tone. Specifically, I should have noticed that Eva's responses to Jonny's questions about her emotional state (e.g., "Do you ever get scared?") were reassuring and empathetic, which suggests that she is not anxious or fearful. I should have also picked up on the fact that Jonny's tone and language become softer and more intimate towards the end of the conversation, which implies that Eva is feeling a sense of comfort and connection with him. } \\
\\ 
\texttt{To improve my strategy, I will focus on analyzing the emotional dynamics of the conversation, paying attention to the tone, language, and body language (if applicable) of the characters. I will also consider the context of the conversation and the relationships between the characters. In this case, I should have scored Anxious: 0, Relieved: 8, Curious: 0, Apathetic: 0.} \\
            }
        }
    }
    & \\ \\

\end{supertabular}
}

\end{document}
